\chapter{Introduction}

 In recent years, the semiconductor industry has experienced significant growth driven by advancements in manufacturing technologies and the increasing demand for highly efficient, reliable, and high-performance integrated circuits. In modern semiconductor fabrication, defect detection and classification play a critical role in maintaining production efficiency and ensuring product reliability. Despite advancements in inspection systems (automatic and semi-automatic), the reduced size of integrated circuits and the increasingly complex nature of defect patterns continue to pose significant challenges to these approaches. Although deep learning techniques have demonstrated promising results in automated defect analysis, their deployment in real-world environments is often constrained by high computational requirements and limited interoperability across platforms. This has led to increasing interest in lightweight, edge-deployable solutions that can operate efficiently in resource-constrained and heterogeneous environments. 


\section{Background of Defect Classes}
\subsection{Semiconductor Manufacturing Overview}
Semiconductor manufacturing is a highly complex, multi-stage process used to fabricate integrated circuits (ICs) on silicon wafers. The manufacturing process involves repeated cycles of lithography, deposition, etching, doping, and polishing, all performed with nanometer-level precision.
Each wafer undergoes numerous processing steps in a clean room environment. Even minor variations in process conditions can introduce defects that adversely affect circuit functionality.


\subsection{What are Wafer Defects and Why They Occur}
Wafer defects refer to unwanted irregularities or anomalies that occur on the surface or within the semiconductor wafer during fabrication. These defects can arise due to:
\begin{itemize}
\item Particle Containment in clean rooms
\item Equipment malfunction or calibration errors
\item Process variation (temperature, pressure, timing)
\item Material imperfections in silicon wafers
\end{itemize}

\subsection{Types of Defects}
\begin{itemize}
     \item \textbf{Structural Defects:}
    These are physical damages to the wafer surface.

    Examples:
    \begin{itemize}
        \item Scratches
        \item Cracks
        \item Breaks in the wafer surface
    \end{itemize}
    Characteristics:
        \begin{itemize}
            \item Often caused by handling or by mechanical stress
            \item Visible in imaging systems
            \item May affect entire circuit functionality
        \end{itemize}
      \item \textbf{Contamination Defects:}
        These are foreign particles or residues on the wafer surface.
        Examples:
    \begin{itemize}
        \item Dust particles
        \item Chemical residues
        \item Metallic contamination
    \end{itemize}
    Characteristics:
        \begin{itemize}
            \item Usually random in distribution
            \item Can be very small but still be impactful
            \item Often introduced during cleaning or pre/post processing steps
        \end{itemize}

        \item \textbf{Pattern-Related Defects:}
        These defects arise when the intended geometric layout of the integrated circuit is incorrectly formed, leading to deviations from the designed pattern.
        Examples:
    \begin{itemize}
        \item Bridging (unwanted connection between circuit lines)
        \item Open circuits (break in expected pattern)
        \item Misalignment of patterns
        \item Missing or extra pattern features
    \end{itemize}
    Characteristics:
        \begin{itemize}
            \item Occur during lithography or etching steps
            \item Directly affect circuit behavior
            \item Harder to detect visually compared to other structural defects
        \end{itemize}
\end{itemize}

\subsection{Traditional Inspection approaches}
\begin{itemize}
     \item \textbf{Optical Inspection Method:}
   This approach uses high-resolution imaging systems to examine semiconductor wafers. The captured images are analyzed to identify irregularities such as scratches, contamination, or pattern deviations during the manufacturing process.
    \subsubsection{How they work:}
     \begin{itemize} 
     \item Capture wafer images
     \item Compare with reference patterns or neighboring dies
     \item Flag deviations as defects
     \end{itemize}
      Limitations:
      \begin{itemize} 
      \item High cost
      \item Sensitive to noise
      \item Struggles with complex patterns in advanced nodes
      \item Requires post-processing interpretation
      \end{itemize}

     \item \textbf{Manual Classification:}  
       Here human experts visually inspect wafer maps or microscope images to identify defect types.
       Limitations:
       \begin{itemize} 
           \item Time-consuming
           \item Subjective (depends on expert experience)
            \item Not scalable for high-volume production
            \item Inconsistent labeling across engineers
       \end{itemize}
\end{itemize}
Due to the limitations of traditional inspection methods, there has been a significant shift toward automated defect classification using machine learning and deep learning techniques. These approaches aim to improve accuracy, scalability, and real-time decision-making in semiconductor manufacturing environments.

\section{Importance of defect classification}
Defect classification is a key component of semiconductor manufacturing, as it influences product quality, overall yield, and the effectiveness of automated production processes. As device geometries shrink and fabrication complexity increases, accurate and efficient defect analysis becomes essential.

\subsection {Quality Improvement}

Defect classification ensures that only wafers meeting stringent quality standards proceed to packaging and deployment. By identifying and categorizing defects (such as scratches, contamination, or pattern irregularities), manufacturers can:
\begin{itemize}
\item Detect faulty dies early in the production cycle
\item Prevent defective chips from reaching customers
\item Maintain reliability in high-performance applications (e.g., automotive, medical, aerospace systems)
\end{itemize}

This is especially crucial in safety-critical domains like automotive embedded systems, where even minor defects can lead to system failures.

\subsection {Yield Enhancement}

Yield refers to the percentage of functional chips produced from a wafer. Defect classification directly contributes to yield improvement by:
\begin{itemize}
\item Identifying recurring defect patterns linked to specific process steps
\item Enabling root cause analysis (e.g., lithography errors, etching issues)
\item Reducing scrap rates and rework costs
\end{itemize} 

By analyzing defect trends over time, manufacturers can optimize fabrication parameters and significantly increase the number of usable chips per wafer, leading to higher profitability.

\subsection {Process Automation}

Traditional defect inspection methods rely heavily on manual review, which is:
\begin{itemize}
\item Time-consuming
\item Subject to human error
\item Not scalable for high-volume production
\end{itemize} 

Automated defect classification systems, especially those based on deep learning, enable:
\begin{itemize}
\item Real-time inspection and classification of wafer defects
\item Consistent and objective decision-making
\item Integration with smart manufacturing pipelines (Industry 4.0)
\end{itemize} 
Automation reduces dependency on human experts and allows fabs to operate efficiently at scale.

\subsection {Cost Reduction}

Accurate defect classification minimizes:
\begin{itemize}
\item Material wastage
\item Rework and repair costs
\item Downtime due to unidentified process issues
\end{itemize}
Early detection and classification prevent defects from propagating through multiple stages, saving significant operational costs.

\subsection {Data-Driven Manufacturing}

Defect classification generates valuable data that can be used for:
\begin{itemize}
\item Predictive maintenance of equipment
\item Statistical process control (SPC)
\item Continuous process improvement
\end{itemize} 

This enables a shift from reactive to proactive manufacturing strategies.
\section {Limitations and Motivation}
\subsection {Limitation with current approach }
\begin{itemize}
\item High-performing deep learning models (e.g.InceptionNet (GoogLeNet), DenseNet ) are computationally expensive
\item Dependence on cloud-based processing, leading to latency and bandwidth issues
\item Limited suitability for real-time defect detection in high-throughput manufacturing
\item Black-box nature of models reduces trust in critical manufacturing decisions
\item Poor generalization across fabrication environments
\item Heavy reliance on large labeled datasets
\item Traditional methods lack robustness, while deep models lack efficiency
\end{itemize} 

\subsection {Motivation for Using Edge AI}
\begin{itemize}
\item Enables real-time defect classification directly on inspection machines
\item Reduces latency by avoiding cloud communication
\item Minimizes network bandwidth usage (important for high-resolution wafer images)
\item Improve data privacy and security (sensitive manufacturing data stays on-site)
\item Supports scalable deployment across multiple inspection units in fabs
\item Allows integration with Industry 4.0 smart manufacturing pipelines
\end{itemize}

In conclusion, as semiconductor manufacturing moves towards fully automated and intelligent systems, the role of efficient and explainable defect classification becomes increasingly important. The convergence of Edge AI and interpretable deep learning offers a promising direction to achieve real-time, reliable, and scalable solutions. This work is driven by the vision of enabling practical, trustworthy edge AI systems that can be seamlessly integrated into next-generation smart manufacturing environments.